% =============================================================================
%  REVISED BODY -- drop-in replacement for the corresponding part of main.tex
%  Preamble, \title/\author block and bibliography are unchanged; this file
%  covers Abstract through Methodology, plus the predictor table.
% =============================================================================

\thispagestyle{firstpage} % Apply custom header/footer for first page
\begin{abstract}
\fontsize{11}{13}\selectfont
\vspace{1em}
\justifying
Low birth weight (LBW), a birth weight below 2.5~kg, is a leading antecedent of
neonatal mortality, developmental delay, and chronic disease in later life.
Conventional analyses dichotomise birth weight at the 2.5~kg threshold,
discarding all information about severity below it, while conventional tree
methods return a partition with no measure of uncertainty. We develop a
Bayesian, tree-based framework that models the full birth-weight distribution
and quantifies risk at the level of an individual covariate profile. Recursive
partitioning is driven by the marginal Dirichlet-multinomial likelihood, whose
conjugacy gives a closed-form splitting criterion and stabilises sparse cells
rather than merely tolerating them. The 2024 U.S. Natality file is condensed
into 672 maternal-infant classes, and birth weight is discretised into deciles
of the \emph{previous} year's LBW distribution, so cutpoints and prior are
estimated on data disjoint from the counts being modelled. Two models are
grown: a full model over all births ($K=11$) and an LBW-only model restricted
to births below the threshold ($K=10$). A two-tier parametric bootstrap
propagates class- and category-level sampling variation through a
10{,}000-tree ensemble. Because the criterion is a likelihood rather than an
impurity measure, the framework supports categorical predictors at a resolution
plain multinomial CART cannot sustain: maternal race enters with seven groups
rather than a binary indicator, its optimal grouping recovered by exhaustive
search over all bipartitions of its levels. Posterior risk profiles are read
off the ensemble with 95\% percentile intervals, and span roughly fivefold
between the extreme profiles.

\vspace{-0.1cm}
\noindent\textbf{keywords}
Risk stratification, low birth weight, Bayesian nonparametrics, classification
and regression trees, Dirichlet-multinomial likelihood

\end{abstract}

\vspace{1cm}
\end{@twocolumnfalse}
]

% =============================================================================
\section{Introduction}
\fontsize{12}{14}\selectfont

Low birth weight (LBW) is defined as a birth weight below 2.5~kg
\citep{lbw_def, kramer1987}. Infants born underweight face substantially
elevated risk of neonatal mortality, developmental delay, and chronic illness
in later life \citep{finch2003}, and the determinants of that risk are
genetic, biological, environmental, and socioeconomic at once. Identifying
which combinations of maternal and infant characteristics carry the greatest
risk is therefore a prerequisite for targeted prevention.

Two limitations recur in the existing literature. First, birth weight is
usually reduced to a binary indicator at the 2.5~kg threshold, which treats a
1.0~kg birth and a 2.4~kg birth as the same event and discards precisely the
gradient that clinical severity follows. Second, methods that are interpretable
enough to guide policy---decision trees in particular---typically report a
single fitted partition with no quantification of how stable that partition is,
while methods that do quantify uncertainty tend to sacrifice the interpretable
structure practitioners rely on.

We address both by pairing recursive partitioning with a Dirichlet-multinomial
(DM) likelihood. Retaining the full discretised birth-weight distribution as a
vector-valued response preserves severity, and placing a conjugate Dirichlet
prior on the category probabilities yields a marginal likelihood available in
closed form, which serves directly as the splitting criterion. The same
conjugacy that makes the criterion tractable also shrinks sparse cells toward
the marginal profile, so the method does not merely tolerate thin subgroups but
is designed for them. This matters practically: it permits a categorical risk
factor to be disaggregated to a resolution a plain multinomial CART could not
support, and it allows a risk profile---a complete covariate combination---to
be read off the fitted ensemble with an interval attached.

% =============================================================================
\section{Literature Review}
\fontsize{12}{14}\selectfont

The landmark meta-analysis of \citet{kramer1987} catalogued the determinants of
LBW, among them maternal anthropometry, inadequate gestational weight gain,
cigarette smoking, malaria infection, and prior adverse pregnancy outcomes. The
factors it identified are strongly interrelated, which makes their marginal
effects difficult to separate and motivates methods that model interactions
directly rather than adjusting for them.

Tree-based methods answer that need. \citet{KITSANTAS2006275} applied CART to
LBW and demonstrated its value for isolating high-risk profiles in
high-dimensional settings, where the fitted partition is itself the
interpretable output. That application, however, inherited both limitations
noted above: the response was the binary LBW indicator, so severity below
2.5~kg was discarded, and the reported tree carried no uncertainty.

A parallel line of work models the full birth-weight distribution conditional
on covariates, using Bayesian nonparametric mixtures \citep{dunson2008} and
copula-based or density-regression techniques \citep{rathjens2023}. These
recover detailed distributional effects but return a fitted density rather than
a rule, and are correspondingly harder to translate into the subgroup
statements that practitioners and policymakers act on. Our contribution sits
between these traditions: the response is distributional, and the estimator is
a tree.

The covariates themselves are well established. Risk is elevated at younger and
older maternal ages, with lower educational attainment, and with limited access
to prenatal care \citep{age_differences_lbw, finch2003, jain2024}, and further
elevated by tobacco use, substance exposure, and air pollution
\citep{standford_med_lbw, lu2020combined}. Incidence differs sharply across
demographic groups: in the United States the LBW rate among Black infants is
roughly double that among White infants, 14.7\% against 7.1\%
\citep{marchofdimes2024}. It is common to encode race as a single binary
contrast on this basis. We do not, for reasons the method makes affordable and
Section~\ref{sec:data} sets out.

% =============================================================================
\section{Data}
\label{sec:data}
\fontsize{12}{14}\selectfont

\paragraph{Source and predictors.}
The analysis uses the 2024 U.S. Natality file compiled by the National Center
for Health Statistics (NCHS), comprising 3.64 million live-birth records; after
complete-case filtering 3.48 million (95.6\%) are retained. Seven maternal and
infant predictors were selected for clinical relevance, generalisability, and
support in the prior literature (Table~\ref{tab:vars}), spanning demographic,
socioeconomic, and biological domains. Their joint level set defines
\[
  I \;=\; \underbrace{2^{5}}_{\text{five binary}}\times
          \underbrace{3}_{\text{marital status}}\times
          \underbrace{7}_{\text{maternal race}} \;=\; 672
\]
mutually exclusive classes, each a distinct maternal-infant profile. Condensing
3.48 million records into 672 rows of counts is what makes an ensemble of
10{,}000 trees tractable while preserving every distinction the predictors draw.

\paragraph{Two predictors are not binary, by design.}
Marital status is reported as unknown for a non-negligible and geographically
clustered share of records; because the public-use file suppresses every
geographic identifier, the mechanism generating that missingness is
unobservable and no conditional imputation is testable. We therefore retain
those records and carry \emph{unknown} as an explicit third level, which
assumes nothing and lets the fitted model report where the unknown group
belongs. Maternal race is carried at its full seven-group resolution rather
than collapsed to a binary contrast or to an ``other'' bucket that would merge
groups with materially different birth-weight profiles. Both choices enlarge
the class space, and both are affordable only because the DM prior stabilises
the sparse cells that result---the point made in Section~\ref{sec:methods}.

\paragraph{Birth-weight categories.}
Birth weight is discretised into $K$ ordered categories. The ten LBW categories
are the deciles of the birth-weight distribution below 2.5~kg, with cutpoints
estimated from the \emph{preceding} year (2023); all births above 2.5~kg form a
single normal-birth-weight (NBW) category. Estimating cutpoints and prior on
2023 while modelling 2024 counts keeps the two disjoint and avoids using the
data twice. Year-to-year proportions are stable enough for the previous year to
be informative: in 2023, 8.7\% of births fell below 2.5~kg. The full model
retains all births ($K=11$); the LBW-only model restricts to births below the
threshold ($K=10$), where the ten categories are conditional on being LBW.
Figure~\ref{fig:prior} shows the resulting prior for both models.

\paragraph{Notation.}
Let $\mathbf{X}\in\prod_{j=1}^{7}\mathcal{L}_j$ be the fixed
$I\times 7$ predictor matrix, $|\mathcal{L}|=(2,3,2,2,2,2,7)$, whose $i$th row
$\mathbf{x}_i$ is one class profile. Let
$\mathbf{Y}=[n_{i,k}]_{i=1,\dots,I}^{\,k=1,\dots,K}$ be the corresponding
$I\times K$ response matrix, so $\mathbf{y}_i=(n_{i,1},\dots,n_{i,K})$ collects
the births of class $i$ across categories and $N_i=\sum_{k}n_{i,k}$ is its
total. Each pair $(\mathbf{x}_i,\mathbf{y}_i)$ thus records how many births of
a given profile fell into each birth-weight category, governed by an unknown
probability vector $\boldsymbol{\theta}_i$.

% =============================================================================
\section{Methodology}
\label{sec:methods}
\fontsize{12}{14}\selectfont

\subsection{The Dirichlet-multinomial model}

Within a class, category counts are multinomial given
$\boldsymbol{\theta}_i$, and a conjugate Dirichlet prior is placed on
$\boldsymbol{\theta}_i$:
\begin{align*}
\mathbf{y}_i \mid \mathbf{x}_i
   &\sim \text{Multinomial}^{\ast}(N_i,\boldsymbol{\theta}_i)\\
\boldsymbol{\theta}_i &\sim \mathrm{Dir}(\boldsymbol{\alpha})\\
\boldsymbol{\theta}_i \mid \mathbf{y}_i,\mathbf{x}_i
   &\sim \mathrm{Dir}(\boldsymbol{\alpha}+\mathbf{y}_i).
\end{align*}
The asterisk denotes omission of the multinomial coefficient: it does not
depend on the parameters of interest, and retaining it would bias the splitting
criterion toward classes with more births. The hyperparameter
$\boldsymbol{\alpha}=(\alpha_1,\dots,\alpha_K)$, with total mass
$\alpha_0=\sum_k\alpha_k$, is set from the previous year's category
proportions, which gives an informative, non-degenerate prior for categories
that may be empty in any given class. We suppress the class index $i$ below;
the derivation is identical for every class.

Writing the Dirichlet density and the outcome model,
\begin{align*}
  p(\boldsymbol{\theta}\mid\boldsymbol{\alpha})
    &= \frac{\Gamma(\alpha_0)}{\prod_{k}\Gamma(\alpha_k)}
       \prod_{k=1}^{K}\theta_k^{\alpha_k-1},
  &
  p(\mathbf{y}\mid\boldsymbol{\theta},\mathbf{x})
    &\propto \prod_{k=1}^{K}\theta_k^{n_k},
\end{align*}
their product is Dirichlet in $\boldsymbol{\theta}$, which exhibits the
conjugacy directly:
\begin{align*}
  p(\mathbf{y},\boldsymbol{\theta}\mid\mathbf{x})
    &\propto \prod_{k=1}^{K}\theta_k^{\,n_k+\alpha_k-1}
     \;\sim\; \mathrm{Dir}(\boldsymbol{\alpha}+\mathbf{y}).
\end{align*}
Integrating out $\boldsymbol{\theta}$ gives the marginal likelihood in closed
form,
\begin{align*}
  p(\mathbf{y}\mid\mathbf{x})
    &= \int p(\mathbf{y},\boldsymbol{\theta}\mid\mathbf{x})\,
       d\boldsymbol{\theta}
     = \frac{\Gamma(\alpha_0)}{\Gamma(N+\alpha_0)}
       \prod_{k=1}^{K}\frac{\Gamma(n_k+\alpha_k)}{\Gamma(\alpha_k)},
\end{align*}
and hence, on the log scale,
\begin{equation}
\label{eq:loglik}
\begin{split}
  \ell(\mathbf{y}\mid\mathbf{x})
    ={}& \log\Gamma(\alpha_0)-\log\Gamma(N+\alpha_0)\\
       &+\sum_{k=1}^{K}\bigl[\log\Gamma(n_k+\alpha_k)-\log\Gamma(\alpha_k)\bigr].
\end{split}
\end{equation}
Because the $\boldsymbol{\theta}_i$ are independent across classes given
$\boldsymbol{\alpha}$, the full marginal likelihood factorises,
\[
  p(\mathbf{Y}\mid\boldsymbol{\alpha},\mathbf{X})
    = \prod_{i=1}^{I}p(\mathbf{y}_i\mid\boldsymbol{\alpha},\mathbf{x}_i),
\]
so on the log scale it is a sum of terms of the form~\eqref{eq:loglik}. This
factorisation is what makes the criterion practical inside a recursive
partitioner: evaluating a candidate split requires summing~\eqref{eq:loglik}
only over the classes routed to the node in question, never over all $I$.

\subsection{Growing the tree}

Trees are grown with the CART algorithm \citep{breiman1984classification},
implemented through a custom split rule in the R package \texttt{rpart}
\citep{rpart_docs}. A candidate split of node $t$ into children $t_L,t_R$ is
scored by the gain in marginal log-likelihood,
\[
  \Delta\ell \;=\; \ell(t_L)+\ell(t_R)-\ell(t),
\]
with $\ell(\cdot)$ the sum of~\eqref{eq:loglik} over the classes at that node.
The split maximising $\Delta\ell$ is taken and the procedure recurses. Replacing
an impurity measure with a likelihood is what carries the prior into the
splitting decision: a candidate that isolates a handful of births is not
rewarded for the spuriously extreme category proportions it produces, because
those proportions are shrunk toward $\boldsymbol{\alpha}$ before being scored.

\subsection{Categorical splits}

A $k$-level unordered predictor admits $2^{k-1}-1$ distinct bipartitions of its
levels, but the conventional strategy evaluates only the $k-1$ splits that are
contiguous in some fixed ordering of those levels. For binary predictors the
distinction is vacuous---one split either way---which is why it is rarely
raised. For maternal race at seven levels it is not: a contiguous search
examines $6$ of the $63$ available groupings, and which six depends on nothing
more than the order in which the levels happen to be declared. We therefore
evaluate all $2^{k-1}-1$ bipartitions under the same $\Delta\ell$ criterion and
select the maximiser. At the level counts in play the cost is negligible, and
without it the seven-group encoding of race would be nominal rather than real.

\subsection{Uncertainty via a two-tier parametric bootstrap}

Uncertainty is propagated through an ensemble of $B=10{,}000$ trees. Each
replicate is generated in two tiers, matching the two ways the observed table
could have come out differently. The first tier redraws how the $T$ total
births distribute across classes,
$\mathbf{n}^{\ast}\sim\text{Multinomial}(T,\mathbf{p})$ with $p_i$ the observed
share of class $i$; the second redraws, for each class, how its $n_i^{\ast}$
births distribute across categories, using the posterior mean
\[
  \hat{\pi}_{i,k}
  =\frac{n_{i,k}+\alpha_k}{N_i+\alpha_0},
\]
held fixed across replicates. A tree is fitted to each resampled table, and
every quantity reported below---category probabilities, aggregate LBW risk,
variable-usage frequencies---is summarised over the ensemble by its mean and
its 95\% percentile interval.

% =============================================================================
%  Predictor table -- replaces the previous "binary predictors" table, which
%  had a duplicated \caption and listed pre-revision variable names.
% =============================================================================
\begin{table}[ht]
\caption{Predictor definitions. Five predictors are binary; marital status and
maternal race are multi-level, giving $2^{5}\cdot3\cdot7=672$ classes.}
\label{tab:vars}
\centering
\small
\begin{tabular}{@{}llp{3.4cm}@{}}
\toprule
\textbf{Field} & \textbf{Levels} & \textbf{Definition} \\ \midrule
\texttt{sex}        & 2 & Male; female \\
\texttt{mager}      & 2 & Maternal age $>33$~yr; $\le 33$~yr \\
\texttt{meduc}      & 2 & High school completed or beyond; otherwise \\
\texttt{precare1st} & 2 & Prenatal care began in the first trimester; otherwise \\
\texttt{cig\_0}     & 2 & Any pre-pregnancy smoking; none \\
\texttt{dmar}       & 3 & Married; unmarried; unknown \\
\texttt{race}       & 7 & White, Black, Hispanic, Asian, AIAN, NHOPI, multiracial \\
\bottomrule
\end{tabular}
\end{table}
